%%fssp_ncor_acrs
L'algorithme et les formules \'etant complexes
seul le principe de calcul est pr\'esent\'e ici.\\
Selon Brenguier-89, l'activit\'e
et le nombre de gouttes mesur\'ees par un FSSP
s'expriment comme 2 fonctions~: \\ 
$$
\left\{
\begin{array}{rcl}
A & = & f_a ( N_c, \tau, \tau_1, \tau_2, \tau_r ) \\
N & = & f_n ( N_c, \tau, \tau_1, \tau_2, \tau_r ) \\
\end{array}
\right.
$$
o\`u~:\\
$N $ : nombre total de gouttes compt\'ees \\
$A $ : Activit\'e \\
$N_c$ : nombre total de gouttes corrig\'e des
co\"{\i}ncidences \\
$\tau $ : temps de transit d'une particule \\
\\
Avec $A=A_m$ et $N=Reset$, le syst\`eme poss\`ede 2
inconnues, $\tau$ et $N_c$. Il est
r\'esolu par it\'erations successives. $N_{acrs}$
et $\tau_{acrs}$ sont les r\'esultats.\\ 
Si $Reset\times f_e \le Crt_0$ alors aucune correction
est calcul\'ee et $N_{acrs}=0$ et $\tau_{acrs}=0$.

